\section{Algorithms and tests}
\label{app:algorithms}

\paragraph{Graph canonicalisation.}
Contraction patterns are encoded as vertex-coloured multigraphs and canonicalised
with \texttt{nauty}~\cite{McKay:2013} through the \texttt{pynauty} bindings. Ports
are modelled by auxiliary vertices so that the five slots of an $F$ are
distinguishable from one another only up to the symmetries the tensor actually
has.

An earlier version used a Weisfeiler--Leman style hash as a fallback when the
bindings were unavailable. It collided: $49$ order-six classes collapsed onto
$39$ keys, undercounting the atlas. The fallback has been removed rather than
fixed. A canonicalisation that is exact at every degree used is a precondition for
every dimension in this paper, and a fallback that is silently approximate is
worse than no fallback.

\paragraph{Tensor contractions.}
Contractions are planned once per invariant and then executed with reduction mod
$p$ after every step, so intermediates never grow. Planning uses a greedy
ordering; \texttt{opt\_einsum}~\cite{Smith:2018} is supported as an optional
accelerator and produces better orderings when present, but it was not installed
in the environment that produced the certified results.

\paragraph{Spinor contractions.}
Spinor-side invariants contract gamma and charge-conjugation matrices along the
closed paths of a port graph. The same reduce-after-every-step discipline applies.

\paragraph{Modular arithmetic.}
Primes are chosen near $2^{15}$ so that a product of two residues fits in a
machine integer with room to accumulate. Reductions are explicit; there is no
reliance on a library's implicit modulus. An overflow guard checks that no
intermediate exceeds the safe range, and fails rather than wrapping.

\paragraph{What the Euler check is for.}
Every Jacobian row is tested against Euler's homogeneity relation. It is fair to
ask what this catches, since homogeneity of the invariants is guaranteed by the
construction. The answer is that the check does not test homogeneity of the
invariant; it tests the \emph{differentiation}. Euler's relation ties the
computed derivative row back to the value of the invariant at the same point, so
a row produced by a mis-amputated contraction --- the realistic failure mode when
derivatives are taken analytically by removing one copy of $F$ --- fails it,
whatever rank that row would otherwise have contributed. A row that fails Euler
is wrong, and this is the only check in the pipeline that says so about a single
row in isolation.

\paragraph{Elimination and determinants.}
Ranks over $\Fp$ come from exact Gaussian elimination. Determinants certifying a
rank are computed twice: by fraction-free Bareiss elimination~\cite{Bareiss:1968},
which keeps every intermediate an integer, and by an independent exact routine.
Agreement is required.

\paragraph{Rational reconstruction.}
CRT combination followed by rational reconstruction~\cite{Wang:1982}, always with
at least one prime held out of the fit for validation. Sampling arguments
elsewhere in the pipeline are Schwartz--Zippel in
character~\cite{Zippel:1979,Schwartz:1980}: a non-zero polynomial cannot vanish at
too many random points, so agreement on independently drawn samples bounds the
probability of coincidence, though it does not eliminate it.

\paragraph{Test structure.}
$\nTraceTests$ tensor-side and $\nBridgeTests$ bridge-side tests. Beyond the
ordinary unit tests, the suite contains three categories that carry more weight:
\begin{itemize}\itemsep2pt
\item \textbf{Controls.} Tests asserting that a quantity does \emph{not} vanish
      where it should not, without which the corresponding vanishing tests are
      vacuous. The $\langle F,F\rangle$ control of appendix~\ref{app:gten} is the
      clearest example.
\item \textbf{Counterfactuals.} Tests that deliberately assume something false
      and assert the consequence --- the G-10 counterfactual, the orientation
      branch flip. These detect an assumption becoming silently unnecessary.
\item \textbf{Negative fixtures.} Objects retained specifically because they
      reproduce a historical error, such as the raw target span.
\end{itemize}

\paragraph{Historical defects, retained as fixtures.}
\begin{enumerate}\itemsep2pt
\item Colliding graph canonicalisation, undercounting the atlas.
\item The raw target span identified with the flow closure, giving
      $\Qten = 0$.
\item An unpinned square-root branch in the frame congruence, silently selecting
      the wrong eigenspace at some primes.
\item A control test calling the projector with a signature flag mistaken for a
      duality flag.
\item An identity check that normalised the two sides of an equation separately.
\item A rank reported from an incomplete evaluation schedule.
\end{enumerate}
Each is now covered by a test that fails if the defect returns. We list them
because the credibility of a computational result rests as much on the errors
that were found and fixed as on the checks that passed.
